\label{sec:background}

\subsection{The Section~203 Coverage Formula}

Section~203(b) of the Voting Rights Act, 52~U.S.C.~\S~10503(b), establishes
the coverage formula. A political subdivision is covered if:

\begin{enumerate}
\item A single language minority group is either:
  \begin{itemize}
  \item More than 5\% of the voting-age citizens of the subdivision, \textit{or}
  \item More than 10,000 voting-age citizens in the subdivision;
  \end{itemize}
  \textit{and}
\item The English literacy rate of the language minority group is below
  the national English literacy rate.
\end{enumerate}

Coverage requires both prongs to be satisfied for the same language
minority group. A jurisdiction with more than 10,000 Spanish-speaking
voting-age citizens but whose English literacy rate is at or above
the national rate is not covered for Spanish. Coverage is language-specific:
a county may be covered for Spanish but not for Vietnamese, or covered for
both, depending on the demographics of each group separately.

\subsection{Coverage Determination Process}

The Census Bureau determines Section~203 coverage using data from the
American Community Survey (ACS) five-year estimates. The Bureau collects
data on English proficiency (self-reported as ``Speak English less than
very well'') and citizenship by county and language group from the ACS
public-use microdata. Coverage determinations are made every ten years,
coinciding with the decennial census cycle.

The most recent determination, effective for redistricting through 2032,
was published at 88 Fed.\ Reg.\ 34,396 (May 30, 2023). This determination
covers 331 jurisdictions (counties, county-equivalents, and townships)
in 31 states. The covered languages include:

\begin{itemize}
\item Spanish (largest coverage --- 232 jurisdictions)
\item Chinese (including Mandarin, Cantonese, and other dialects)
\item Filipino (Tagalog)
\item Vietnamese
\item Korean
\item Japanese
\item Khmer (Cambodian)
\item Numerous Native American languages, including Navajo, Yupik,
  Apache, Cherokee, Choctaw, and others
\item Alaska Native languages, including Inupiat and Yupik dialects
\end{itemize}

\subsection{Covered States and Jurisdictions}

The states with the largest number of Section~203 covered jurisdictions
under the 2023 determination are:

\begin{table}[ht]
\centering
\caption{States with Largest Section~203 Coverage (2023 Determination)}
\label{tab:coverage}
\begin{tabular}{lll}
\toprule
State & Jurisdictions Covered & Primary Languages \\
\midrule
Texas & 67 & Spanish, Vietnamese \\
California & 29 & Spanish, Chinese, Vietnamese, Korean, Filipino \\
Arizona & 21 & Spanish, Navajo, Apache, Yaqui \\
New Mexico & 18 & Spanish, Navajo \\
New York & 14 & Spanish, Chinese, Korean \\
Florida & 10 & Spanish, Haitian Creole \\
Alaska & 9 & Alaska Native languages \\
\bottomrule
\end{tabular}
\begin{tablenotes}
\small
\item Source: 88 Fed.\ Reg.\ 34,396 (2023). ``Jurisdictions'' include
counties, county-equivalents, and townships. Alaska is covered based
on Native village areas.
\end{tablenotes}
\end{table}

\subsection{What Section~203 Requires}

Once a jurisdiction is covered, Section~203 requires it to provide all
election materials in both English and the covered language. This includes:
voter registration forms; voting instructions posted at polling places;
absentee ballot materials; and, where applicable, sample ballots and
ballot measures. The requirement runs to the county election administration,
regardless of how the county is divided among congressional districts.

Enforcement is through the Department of Justice, which has authority to
bring civil actions under 52~U.S.C.~\S~10308 to enforce Section~203.
Private plaintiffs may also bring Section~203 enforcement actions.\footnote{
See \textit{Chinatown Voter Education Alliance v.\ Ravitz}, 2009 WL 6445161
(S.D.N.Y.\ 2009) (Section~203 enforcement action by private plaintiffs
regarding Chinese-language materials in New York City).}

\subsection{Section~203 and Redistricting: The Administrative Link}

Section~203 is not itself a redistricting requirement: it does not mandate
how district lines are drawn, and it does not create a right to
district-level representation for language minority groups (Section~2
does that). But redistricting affects Section~203 administration in
two ways.

First, county splits increase administrative burden. When a
Section~203-covered county is split between two congressional districts,
the county election office must manage district-level materials (ballots,
candidate information) for two separate federal elections simultaneously,
in the required language(s) for each district. While this is legally
manageable, it increases cost and error risk.

Second, community integrity matters for effective language assistance.
Language minority communities that are geographically compact---an
urban Chinatown, a Rio Grande Valley colonia, a Navajo chapter
territory---are best served by remaining within a single congressional
district, where a single member of Congress is accountable to the community
and where community organizations can coordinate outreach. When such a
community is split across district lines, the members of Congress
representing each fragment have reduced incentives to serve the community's
interests, and outreach organizations face a more complex advocacy environment.
